\documentclass[11pt]{article}

% pdfLaTeX-friendly (no fontspec)
\usepackage[margin=1in]{geometry}
\usepackage[T1]{fontenc}
\usepackage[utf8]{inputenc}
\usepackage{lmodern}

\usepackage{amsmath,amssymb}
\usepackage{xcolor}
\usepackage{hyperref}
\usepackage{braket} % provides \ket{}, \bra{}, \braket{}
\usepackage{listings}

% Lean listings setup (handles Lean Unicode in code blocks)
\lstdefinelanguage{Lean}{
  morekeywords={import,open,scoped,namespace,section,end,variable,def,abbrev,lemma,theorem,by,fun,have,calc,refine,intro,exact,classical,simp,funext,by_cases,omit,structure,noncomputable,unfold,rcases,obtain,linarith,aesop,grind,where},
  sensitive=true,
  morecomment=[l]{--},
  morecomment=[s]{/-}{-/},
  morestring=[b]"
}

\lstset{
  language=Lean,
  basicstyle=\ttfamily\footnotesize,
  columns=fullflexible,
  breaklines=true,
  frame=single,
  rulecolor=\color{black!25},
  keywordstyle=\color{blue!70!black},
  commentstyle=\color{green!40!black},
  stringstyle=\color{orange!60!black},
  showstringspaces=false,
  keepspaces=true,
  literate=
    {ℕ}{{\ensuremath{\mathbb{N}}}}1
    {ℚ}{{\ensuremath{\mathbb{Q}}}}1
    {ℝ}{{\ensuremath{\mathbb{R}}}}1
    {𝕜}{{\ensuremath{\mathbb{k}}}}1
    {→}{{\ensuremath{\to}}}1
    {↔}{{\ensuremath{\leftrightarrow}}}1
    {∑}{{\ensuremath{\sum}}}1
    {⊆}{{\ensuremath{\subseteq}}}1
    {∈}{{\ensuremath{\in}}}1
    {∀}{{\ensuremath{\forall}}}1
    {∃}{{\ensuremath{\exists}}}1
    {≠}{{\ensuremath{\neq}}}1
    {↑}{{\ensuremath{\uparrow}}}1
    {·}{{\ensuremath{\bullet}}}1
    {⦃}{{\{} }1
    {⦄}{{\}} }1
    {“}{{``}}1
    {”}{{''}}1
    {≤}{{\ensuremath{\le}}}1
    {≥}{{\ensuremath{\ge}}}1
}

\title{Analytic Solution of a Distance-2 SSLP Linear Program and Lean Verification (Worked $((5,2,2))$ Example)}
\author{}
\date{}

\begin{document}
\maketitle
\tableofcontents

\section{Big picture}

In the SSLP pipeline, once we enforce the \emph{classical union distance} $d(C)=2$ for the support
$C=C_{S_0}(\mathbf w)\cup C_{S_1}(\mathbf w)$,
all single-qubit $X_i$ and $Y_i$ Knill--Laflamme constraints vanish automatically.
The remaining weight-1 Knill--Laflamme conditions reduce to a system of \emph{linear equalities}
coming from the diagonal operators $Z_i$ (``$Z$-marginal equalities'').

This note has four goals:
\begin{itemize}
\item clearly state the resulting LP feasibility problem (in mathematical terms),
\item solve the LP \emph{by hand} and derive the full analytic (general) solution,
\item explain how Lean can \emph{verify} the claimed general solution without being asked to ``search'' for one,
\item record the exact Aristotle prompt and explain the generated Lean verifier line-by-line.
\end{itemize}

\section{The mathematical problem (LP feasibility) for the worked $((5,2,2))$ example}

\subsection{Parameters and supports}

We consider the concrete $((5,2,2))$ SSLP example with
\[
n=5,\qquad m=7,\qquad \mathbf w=(1,1,2,2,2),\qquad (S_0,S_1)=(0,4).
\]
Let
\[
C_S(\mathbf w):=\{x\in\{0,1\}^5:\ \langle\mathbf w,x\rangle\equiv S \!\!\!\pmod 7\},
\qquad
\langle\mathbf w,x\rangle:=\sum_{i=1}^5 w_i x_i \pmod 7.
\]
The two residue classes used in the worked example are
\[
C_0=\{x_0,x_1,x_2\}=\{00000,\,01111,\,10111\},
\qquad
C_1=\{y_0,\dots,y_5\}=\{00011,\,00101,\,00110,\,11001,\,11010,\,11100\}.
\]
We work in the design regime where the classical union $C:=C_0\cup C_1$ has $d(C)=2$ (no Hamming-$1$ neighbors).
Then for every site $i$ and logical labels $j,k$,
\[
\bra{j_L}X_i\ket{k_L}=\bra{j_L}Y_i\ket{k_L}=0,
\]
so only the diagonal $Z_i$ constraints remain.

\subsection{Variables (probabilities)}

Introduce probability variables for the two logical states:
\[
a_j := p_{0,x_j}\ \ (j=0,1,2),\qquad b_k := p_{1,y_k}\ \ (k=0,1,2,3,4,5).
\]
They satisfy
\[
a_0+a_1+a_2=1,\qquad b_0+b_1+b_2+b_3+b_4+b_5=1,
\]
and the nonnegativity constraints
\[
a_0,a_1,a_2\ge 0,\qquad b_0,\dots,b_5\ge 0.
\]

\subsection{Linear constraints from $Z$-marginals}

Because $Z_i\ket{x}=(-1)^{x_i}\ket{x}$, for any state supported on basis strings
the expectation $\bra{\psi}Z_i\ket{\psi}$ is a linear combination of the probabilities.
The remaining Knill--Laflamme constraints are, for each site $i=1,\dots,5$,
\[
\bra{0_L}Z_i\ket{0_L}=\bra{1_L}Z_i\ket{1_L}.
\]
For the concrete supports $C_0,C_1$ above, this yields the system:
\begin{align}
\label{eq:LP_i1}
(i = 1)\quad a_{0} + a_{1} - a_{2}
&= b_{0} + b_{1} + b_{2} - b_{3} - b_{4} - b_{5}, \\
\label{eq:LP_i2}
(i = 2)\quad a_{0} - a_{1} + a_{2}
&= b_{0} + b_{1} + b_{2} - b_{3} - b_{4} - b_{5}, \\
\label{eq:LP_i3}
(i = 3)\quad a_{0} - a_{1} - a_{2}
&= b_{0} - b_{1} - b_{2} + b_{3} + b_{4} - b_{5}, \\
\label{eq:LP_i4}
(i = 4)\quad a_{0} - a_{1} - a_{2}
&= -b_{0} + b_{1} - b_{2} + b_{3} - b_{4} + b_{5}, \\
\label{eq:LP_i5}
(i = 5)\quad a_{0} - a_{1} - a_{2}
&= -b_{0} - b_{1} + b_{2} - b_{3} + b_{4} + b_{5}.
\end{align}

\subsection{LP as a feasibility problem}

Over a linear ordered field $\mathbb{k}$ (typically $\mathbb{Q}$ or $\mathbb{R}$), define the feasible set
\[
\mathcal F
:=
\Bigl\{
(a_0,a_1,a_2,b_0,\dots,b_5)\in\mathbb{k}^9:
\ \text{\eqref{eq:LP_i1}--\eqref{eq:LP_i5} hold, both normalizations hold, and all variables are }\ge 0
\Bigr\}.
\]
Equivalently, this is the standard LP form ``find $p$ such that $A p = c$ and $p\ge 0$'' (with objective $\min 0$).

\section{Solving the LP equalities by hand (analytic general solution)}

We now derive the full family of solutions in closed form, and then impose inequalities to obtain the parameter region.

\subsection{Step A: pin the $a$-variables}

Subtract \eqref{eq:LP_i2} from \eqref{eq:LP_i1}; the right-hand sides are equal, hence
\[
(a_0+a_1-a_2)-(a_0-a_1+a_2)=0 \quad\Longrightarrow\quad a_1=a_2.
\]
With $a_0+a_1+a_2=1$, we obtain
\[
a_1=a_2=\frac{1-a_0}{2},
\qquad
a_0-a_1-a_2=a_0-(1-a_0)=2a_0-1.
\]

\subsection{Step B: relations among the $b$-variables}

The left-hand sides of \eqref{eq:LP_i3}, \eqref{eq:LP_i4}, \eqref{eq:LP_i5} are identical. Subtract them pairwise:
\[
\eqref{eq:LP_i3}-\eqref{eq:LP_i4}:\quad b_0-b_1+b_4-b_5=0,
\]
\[
\eqref{eq:LP_i3}-\eqref{eq:LP_i5}:\quad b_0-b_2+b_3-b_5=0,
\]
\[
\eqref{eq:LP_i4}-\eqref{eq:LP_i5}:\quad b_1-b_2+b_3-b_4=0.
\]
Define free parameters
\[
u:=b_4,\qquad v:=b_5.
\]
From the first relation we can already write
\[
b_0=b_1+v-u. \tag{B1}
\]

\subsection{Step C: solve for all $b$'s and pin $a_0$}

Use \eqref{eq:LP_i1}. Since $a_1=a_2$, the left-hand side becomes $a_0$, hence
\[
a_0=b_0+b_1+b_2-b_3-u-v. \tag{C1}
\]

Next, use \eqref{eq:LP_i3}. Its left-hand side is $a_0-a_1-a_2=2a_0-1$, hence
\[
2a_0-1=b_0-b_1-b_2+b_3+u-v. \tag{C2}
\]

Now use the third relation from Step B to express $b_3$:
\[
b_1-b_2+b_3-u=0\quad\Longrightarrow\quad b_3=b_2+u-b_1. \tag{C3}
\]

Substitute (B1) and (C3) into (C1):
\[
\begin{aligned}
a_0
&=(b_1+v-u)+b_1+b_2-(b_2+u-b_1)-u-v\\
&=3(b_1-u). \tag{C4}
\end{aligned}
\]
Similarly, substitute (B1) and (C3) into (C2):
\[
\begin{aligned}
2a_0-1
&=(b_1+v-u)-b_1-b_2+(b_2+u-b_1)+u-v\\
&=-b_1+u. \tag{C5}
\end{aligned}
\]
Combine (C4) and (C5):
\[
2\cdot 3(b_1-u)-1=-b_1+u
\quad\Longrightarrow\quad
7(b_1-u)=1
\quad\Longrightarrow\quad
b_1=u+\frac{1}{7}. \tag{C6}
\]
Then (C4) gives
\[
a_0=3\cdot\frac{1}{7}=\frac{3}{7},
\qquad
a_1=a_2=\frac{1-a_0}{2}=\frac{2}{7}. \tag{C7}
\]

Now solve for the remaining $b$'s using normalization and (C3).
First, from (B1) and (C6),
\[
b_0=b_1+v-u=\frac{1}{7}+v. \tag{C8}
\]
From (C3) and (C6),
\[
b_3=b_2+u-\Bigl(u+\frac{1}{7}\Bigr)=b_2-\frac{1}{7}. \tag{C9}
\]
Apply the $b$-normalization:
\[
b_0+b_1+b_2+b_3+u+v=1.
\]
Substitute (C6), (C8), (C9):
\[
\left(\frac{1}{7}+v\right)+\left(\frac{1}{7}+u\right)+b_2+\left(b_2-\frac{1}{7}\right)+u+v=1,
\]
which simplifies to
\[
2b_2+2u+2v=\frac{6}{7}
\quad\Longrightarrow\quad
b_2=\frac{3}{7}-u-v. \tag{C10}
\]
Then (C9) yields
\[
b_3=\frac{2}{7}-u-v. \tag{C11}
\]

\subsection{Step D: impose inequalities (parameter region)}

Nonnegativity implies
\[
u=b_4\ge 0,\qquad v=b_5\ge 0,
\qquad
b_3=\frac{2}{7}-u-v\ge 0\ \Longrightarrow\ u+v\le \frac{2}{7}.
\]
This automatically implies $b_2=\frac{3}{7}-u-v\ge 0$.
Therefore the admissible region is the triangle
\[
\boxed{u\ge 0,\qquad v\ge 0,\qquad u+v\le \frac{2}{7}.}
\]

\subsection{General solution (final form)}

Collecting (C7), (C8), (C10), (C11) and the definitions $b_4=u$, $b_5=v$, we obtain the full analytic solution:
\[
\boxed{
\begin{aligned}
a_0&=\frac{3}{7},\qquad a_1=a_2=\frac{2}{7},\\
b_0&=\frac{1}{7}+v,\qquad b_1=\frac{1}{7}+u,\qquad
b_2=\frac{3}{7}-u-v,\qquad
b_3=\frac{2}{7}-u-v,\qquad
b_4=u,\qquad b_5=v,
\end{aligned}}
\qquad
\boxed{u\ge 0,\ v\ge 0,\ u+v\le\frac{2}{7}.}
\]

\section{How Lean verifies the solution without ``solving the LP again''}

\subsection{What Lean proves}

Lean is asked to check that the proposed parametrization is \emph{exactly} the feasible set.
This is the same as proving two logical implications (two containments):

\paragraph{Soundness.}
\[
\forall u v,\ \mathrm{Region}(u,v)\Rightarrow \mathrm{Feasible}(\mathrm{sol}(u,v)).
\]
This is pure substitution: unfold definitions and verify all constraints.

\paragraph{Completeness.}
\[
\forall p,\ \mathrm{Feasible}(p)\Rightarrow \mathrm{Region}(p.b4,p.b5)\wedge p=\mathrm{sol}(p.b4,p.b5).
\]
This is the ``hard direction,'' but crucially it involves \emph{no parameter search}:
we simply set $u:=p.b4$ and $v:=p.b5$ and prove the remaining coordinates are forced by the linear equalities.

\paragraph{Uniqueness (optional).}
\[
\mathrm{sol}(u,v)=\mathrm{sol}(u',v') \Rightarrow u=u'\wedge v=v'.
\]
This is immediate since the structure fields include $b_4=u$ and $b_5=v$.

\subsection{Why this scales better than asking Lean to ``solve''}

Lean never runs an LP algorithm. It checks a human-guided proof:
\begin{itemize}
\item derive $a_1=a_2$ from two equalities,
\item derive the three $b$-relations by subtracting equalities,
\item derive the closed form by combining a small number of equalities and normalizations,
\item derive the region constraints from nonnegativity.
\end{itemize}
For this example, tactics like \texttt{linarith} and \texttt{ring\_nf} are sufficient.
Automation (e.g.\ \texttt{grind}) is used as a convenience to discharge straightforward linear arithmetic goals.

\section{Prompt sent to Aristotle (specification)}

\noindent The exact prompt sent to Aristotle is recorded here for reproducibility.

\begin{verbatim}
% =====================================================================
% Specification for Lean/Aristotle: verifying an analytic general solution
% for the SSLP linear (LP) stage in the concrete ((5,2,2)) example (m=7).
% =====================================================================

\section{LP verification task for the worked $((5,2,2))$ example (distance $2$ case)}

\subsection{What the problem is}

We consider the concrete $((5,2,2))$ SSLP example with
\[
n=5,\qquad m=7,\qquad \mathbf w=(1,1,2,2,2),\qquad (S_0,S_1)=(0,4),
\]
and the two subset-sum residue supports
\[
C_0=\{x_0,x_1,x_2\}=\{00000,\,01111,\,10111\},\qquad
C_1=\{y_0,\dots,y_5\}=\{00011,\,00101,\,00110,\,11001,\,11010,\,11100\}.
\]
We are in the design regime where the classical union
\(
C:=C_0\cup C_1
\)
satisfies $d(C)=2$ (no Hamming-$1$ neighbors inside $C$). Consequently, the $X_i/Y_i$ Knill--Laflamme constraints
vanish automatically, and the SSLP linear stage reduces to the \emph{single-qubit $Z_i$ marginal equalities}.

\paragraph{Variables (probabilities).}
Introduce probability variables (nonnegative reals) for the two logical states:
\[
a_j := p_{0,x_j}\ \ (j=0,1,2),\qquad
b_k := p_{1,y_k}\ \ (k=0,1,2,3,4,5).
\]
These must satisfy the normalizations
\[
a_0+a_1+a_2=1,\qquad b_0+b_1+b_2+b_3+b_4+b_5=1,
\]
and nonnegativity
\[
a_0,a_1,a_2\ge 0,\qquad b_0,\dots,b_5\ge 0.
\]

\paragraph{The linear (LP) constraints.}
For each site $i\in\{1,2,3,4,5\}$, the KL condition for $Z_i$ is
\[
\bra{0_L}Z_i\ket{0_L}=\bra{1_L}Z_i\ket{1_L}.
\]
Since $Z_i\ket{x}=(-1)^{x_i}\ket{x}$, these are \emph{linear equalities in the probabilities}.
For the concrete supports above, these equalities are exactly the system:
\begin{align}
\label{eq:LP_i1}
(i = 1)\quad a_{0} + a_{1} - a_{2}
&= b_{0} + b_{1} + b_{2} - b_{3} - b_{4} - b_{5}, \\
\label{eq:LP_i2}
(i = 2)\quad a_{0} - a_{1} + a_{2}
&= b_{0} + b_{1} + b_{2} - b_{3} - b_{4} - b_{5}, \\
\label{eq:LP_i3}
(i = 3)\quad a_{0} - a_{1} - a_{2}
&= b_{0} - b_{1} - b_{2} + b_{3} + b_{4} - b_{5}, \\
\label{eq:LP_i4}
(i = 4)\quad a_{0} - a_{1} - a_{2}
&= -b_{0} + b_{1} - b_{2} + b_{3} - b_{4} + b_{5}, \\
\label{eq:LP_i5}
(i = 5)\quad a_{0} - a_{1} - a_{2}
&= -b_{0} - b_{1} + b_{2} - b_{3} + b_{4} + b_{5}.
\end{align}

\paragraph{LP as a feasibility problem.}
Over a linear ordered field $\mathbb{k}$ (typically $\mathbb{Q}$ or $\mathbb{R}$), the LP feasibility set is
\[
\mathcal F
:=
\Bigl\{
(a_0,a_1,a_2,b_0,\dots,b_5)\in\mathbb{k}^9:
\ \text{\eqref{eq:LP_i1}--\eqref{eq:LP_i5} hold, normalizations hold, and all variables are }\ge 0
\Bigr\}.
\]
Equivalently, this is the standard LP form ``find $p$ such that $A p = c$ and $p\ge 0$''.

\subsection{What the solution is (claimed analytic general solution)}

Define free parameters
\[
u := b_4,\qquad v := b_5.
\]
The claimed \emph{general solution} of the LP feasibility set $\mathcal F$ is the following affine parametrization:
\begin{equation}
\label{eq:claimed-solution}
\boxed{
\begin{aligned}
a_0&=\frac{3}{7},\qquad a_1=a_2=\frac{2}{7},\\
b_0&=\frac{1}{7}+v,\qquad b_1=\frac{1}{7}+u,\qquad
b_2=\frac{3}{7}-u-v,\qquad
b_3=\frac{2}{7}-u-v,\qquad
b_4=u,\qquad b_5=v.
\end{aligned}}
\end{equation}
The admissible parameter region is the triangle
\begin{equation}
\label{eq:region}
\boxed{
u\ge 0,\qquad v\ge 0,\qquad u+v\le \frac{2}{7}.
}
\end{equation}
(Indeed, $b_3\ge 0$ forces $u+v\le 2/7$, and this automatically implies $b_2\ge 0$.)

Thus the claim is:
\[
\mathcal F
=
\Bigl\{
\mathrm{sol}(u,v)\ :\ (u,v)\in\mathbb{k}^2\ \text{satisfying \eqref{eq:region}}
\Bigr\},
\]
where $\mathrm{sol}(u,v)$ is the $9$-tuple given by \eqref{eq:claimed-solution}.

\subsection{What we want Lean to check (prove)}

Lean is \emph{not} asked to solve the LP automatically.
Instead, Lean is asked to \emph{verify} that the proposed parametrization
\eqref{eq:claimed-solution}--\eqref{eq:region} is exactly the feasible set.

% End of specification.
\end{verbatim}

\section{Explanation of the Lean verifier (what it does step by step)}

This section explains the generated Lean file in the same ``code block + meaning'' style as typical project notes.

\subsection{Imports and environment}

\begin{lstlisting}
import Mathlib
noncomputable section
\end{lstlisting}

\begin{itemize}
\item \texttt{Mathlib} provides ordered fields (\texttt{ℚ}, \texttt{ℝ}) and automation tactics used by Aristotle.
\item \texttt{noncomputable} signals that we are not trying to extract an algorithm; we are proving propositions.
\end{itemize}

\subsection{The variables as a record \texttt{Vars}}

\begin{lstlisting}
structure Vars (k : Type*) where
  a0 : k
  a1 : k
  a2 : k
  b0 : k
  b1 : k
  b2 : k
  b3 : k
  b4 : k
  b5 : k
\end{lstlisting}

This is the mathematical vector $(a_0,a_1,a_2,b_0,\dots,b_5)\in k^9$, but with named fields so constraints remain readable.

\subsection{Encoding feasibility: the predicate \texttt{Feasible}}

\begin{lstlisting}
def Feasible {k : Type*} [Field k] [LinearOrder k] [IsStrictOrderedRing k] (p : Vars k) : Prop :=
  p.a0 ≥ 0 ∧ p.a1 ≥ 0 ∧ p.a2 ≥ 0 ∧
  p.b0 ≥ 0 ∧ p.b1 ≥ 0 ∧ p.b2 ≥ 0 ∧ p.b3 ≥ 0 ∧ p.b4 ≥ 0 ∧ p.b5 ≥ 0 ∧
  p.a0 + p.a1 + p.a2 = 1 ∧
  p.b0 + p.b1 + p.b2 + p.b3 + p.b4 + p.b5 = 1 ∧
  p.a0 + p.a1 - p.a2 = p.b0 + p.b1 + p.b2 - p.b3 - p.b4 - p.b5 ∧
  p.a0 - p.a1 + p.a2 = p.b0 + p.b1 + p.b2 - p.b3 - p.b4 - p.b5 ∧
  p.a0 - p.a1 - p.a2 = p.b0 - p.b1 - p.b2 + p.b3 + p.b4 - p.b5 ∧
  p.a0 - p.a1 - p.a2 = -p.b0 + p.b1 - p.b2 + p.b3 - p.b4 + p.b5 ∧
  p.a0 - p.a1 - p.a2 = -p.b0 - p.b1 + p.b2 - p.b3 + p.b4 + p.b5
\end{lstlisting}

\begin{itemize}
\item The first block encodes nonnegativity.
\item The next two equalities are the two normalizations.
\item The final five equalities are exactly \eqref{eq:LP_i1}--\eqref{eq:LP_i5}.
\item The assumptions \texttt{[Field k] [LinearOrder k] [IsStrictOrderedRing k]} ensure Lean has an ordered field so it can reason about inequalities and use \texttt{linarith}.
\end{itemize}

\subsection{The parameter region and explicit solution family}

\begin{lstlisting}
def Region {k : Type*} [Field k] [LinearOrder k] [IsStrictOrderedRing k] (u v : k) : Prop :=
  u ≥ 0 ∧ v ≥ 0 ∧ u + v ≤ 2/7

def sol {k : Type*} [Field k] [LinearOrder k] [IsStrictOrderedRing k] (u v : k) : Vars k :=
  { a0 := 3/7, a1 := 2/7, a2 := 2/7,
    b0 := 1/7 + v, b1 := 1/7 + u, b2 := 3/7 - u - v,
    b3 := 2/7 - u - v, b4 := u, b5 := v }
\end{lstlisting}

This encodes exactly the analytic general solution and the triangular region.

\subsection{Soundness: \texttt{Region u v → Feasible (sol u v)}}

\begin{lstlisting}
theorem soundness (u v : k) :
  Region u v → Feasible (sol u v) := by
    unfold Region Feasible sol
    grind
\end{lstlisting}

After unfolding, the goal is a list of explicit linear inequalities/equalities, and \texttt{grind} closes them.
Mathematically, this is ``plug in the formulas and check the constraints.''

\subsection{Completeness: proving no feasible point is missed}

The Lean file mirrors the hand proof by providing intermediate lemmas, each corresponding to one logical step.

\begin{itemize}
\item \texttt{completeness\_step\_A}: derives $a_1=a_2$ and related rewrites from the first two equalities + normalization.
\item \texttt{completeness\_step\_B}: derives the three $b$-relations by subtracting the last three equations.
\item \texttt{completeness\_step\_C}: pins all coordinates in terms of $p.b4$ and $p.b5$ (i.e.\ $u:=b_4$, $v:=b_5$).
\item \texttt{completeness\_step\_D}: derives the region constraints from nonnegativity and the formula for $b_3$.
\end{itemize}

The main completeness theorem then packages these into:
\[
\mathrm{Region}(p.b4,p.b5)\ \wedge\ p=\mathrm{sol}(p.b4,p.b5).
\]

\subsection{Uniqueness: parameter injectivity}

\begin{lstlisting}
theorem uniqueness (u v u' v' : k) :
  Region u v → Region u' v' → sol u v = sol u' v' → u = u' ∧ v = v' := by
  intro h₁ h₂ h₃
  simp_all +decide [ sol ]
\end{lstlisting}

Since $b_4=u$ and $b_5=v$ are stored directly, equality of records forces equality of parameters.

\appendix
\section{Lean verifier file (full listing)}

\begin{lstlisting}
/-
This file was generated by Aristotle.

We verify the analytic general solution for the SSLP linear (LP) stage in the concrete ((5,2,2)) example.
-/

import Mathlib

set_option linter.mathlibStandardSet false

open scoped BigOperators
open scoped Real
open scoped Nat
open scoped Classical
open scoped Pointwise

set_option maxHeartbeats 0
set_option maxRecDepth 4000
set_option synthInstance.maxHeartbeats 20000
set_option synthInstance.maxSize 128

set_option relaxedAutoImplicit false
set_option autoImplicit false

noncomputable section

structure Vars (k : Type*) where
  a0 : k
  a1 : k
  a2 : k
  b0 : k
  b1 : k
  b2 : k
  b3 : k
  b4 : k
  b5 : k

def Feasible {k : Type*} [Field k] [LinearOrder k] [IsStrictOrderedRing k] (p : Vars k) : Prop :=
  p.a0 ≥ 0 ∧ p.a1 ≥ 0 ∧ p.a2 ≥ 0 ∧
  p.b0 ≥ 0 ∧ p.b1 ≥ 0 ∧ p.b2 ≥ 0 ∧ p.b3 ≥ 0 ∧ p.b4 ≥ 0 ∧ p.b5 ≥ 0 ∧
  p.a0 + p.a1 + p.a2 = 1 ∧
  p.b0 + p.b1 + p.b2 + p.b3 + p.b4 + p.b5 = 1 ∧
  p.a0 + p.a1 - p.a2 = p.b0 + p.b1 + p.b2 - p.b3 - p.b4 - p.b5 ∧
  p.a0 - p.a1 + p.a2 = p.b0 + p.b1 + p.b2 - p.b3 - p.b4 - p.b5 ∧
  p.a0 - p.a1 - p.a2 = p.b0 - p.b1 - p.b2 + p.b3 + p.b4 - p.b5 ∧
  p.a0 - p.a1 - p.a2 = -p.b0 + p.b1 - p.b2 + p.b3 - p.b4 + p.b5 ∧
  p.a0 - p.a1 - p.a2 = -p.b0 - p.b1 + p.b2 - p.b3 + p.b4 + p.b5

def Region {k : Type*} [Field k] [LinearOrder k] [IsStrictOrderedRing k] (u v : k) : Prop :=
  u ≥ 0 ∧ v ≥ 0 ∧ u + v ≤ 2/7

def sol {k : Type*} [Field k] [LinearOrder k] [IsStrictOrderedRing k] (u v : k) : Vars k :=
  { a0 := 3/7, a1 := 2/7, a2 := 2/7,
    b0 := 1/7 + v, b1 := 1/7 + u, b2 := 3/7 - u - v,
    b3 := 2/7 - u - v, b4 := u, b5 := v }

theorem soundness {k : Type*} [Field k] [LinearOrder k] [IsStrictOrderedRing k] (u v : k) :
  Region u v → Feasible (sol u v) := by
    unfold Region Feasible sol;
    grind

lemma completeness_step_A {k : Type*} [Field k] [LinearOrder k] [IsStrictOrderedRing k] (p : Vars k) (h : Feasible p) :
  p.a1 = p.a2 ∧ p.a1 = (1 - p.a0) / 2 ∧ p.a0 - p.a1 - p.a2 = 2 * p.a0 - 1 := by
    obtain ⟨ ha1, ha2, ha3, ha4, hb ⟩ := h;
    grind

lemma completeness_step_B {k : Type*} [Field k] [LinearOrder k] [IsStrictOrderedRing k] (p : Vars k) (h : Feasible p) :
  p.b0 - p.b1 + p.b4 - p.b5 = 0 ∧
  p.b0 - p.b2 + p.b3 - p.b5 = 0 ∧
  p.b1 - p.b2 + p.b3 - p.b4 = 0 := by
    exact ⟨ by linarith [ h.2.2.2.2.2.2.2.2.2.2 ], by linarith [ h.2.2.2.2.2.2.2.2.2.2 ], by linarith [ h.2.2.2.2.2.2.2.2.2.2 ] ⟩

lemma completeness_step_C {k : Type*} [Field k] [LinearOrder k] [IsStrictOrderedRing k] (p : Vars k) (h : Feasible p) :
  p.a0 = 3/7 ∧ p.a1 = 2/7 ∧ p.a2 = 2/7 ∧
  p.b0 = 1/7 + p.b5 ∧ p.b1 = 1/7 + p.b4 ∧
  p.b2 = 3/7 - p.b4 - p.b5 ∧ p.b3 = 2/7 - p.b4 - p.b5 := by
    obtain ⟨h_a0, h_a1, h_a2, h_b0, h_b1, h_b2, h_b3, h_b4, h_b5, h_sum_a, h_sum_b, h_eq1, h_eq2, h_eq3, h_eq4, h_eq5⟩ := h;
    grind

lemma completeness_step_D {k : Type*} [Field k] [LinearOrder k] [IsStrictOrderedRing k] (p : Vars k) (h : Feasible p) :
  Region p.b4 p.b5 := by
    have := completeness_step_C p h;
    exact ⟨ h.2.2.2.2.2.2.2.1, h.2.2.2.2.2.2.2.2.1, by linarith [ h.2.2.2.2.2.2.1 ] ⟩

theorem completeness {k : Type*} [Field k] [LinearOrder k] [IsStrictOrderedRing k] (p : Vars k) (h : Feasible p) :
  Region p.b4 p.b5 ∧ p = sol p.b4 p.b5 := by
    unfold sol;
    refine' ⟨ completeness_step_D p h, _ ⟩;
    have h_eq :
      p.a0 = 3 / 7 ∧ p.a1 = 2 / 7 ∧ p.a2 = 2 / 7 ∧
      p.b0 = 1 / 7 + p.b5 ∧ p.b1 = 1 / 7 + p.b4 ∧
      p.b2 = 3 / 7 - p.b4 - p.b5 ∧ p.b3 = 2 / 7 - p.b4 - p.b5 := by
      exact?;
    cases p
    aesop

theorem uniqueness {k : Type*} [Field k] [LinearOrder k] [IsStrictOrderedRing k] (u v u' v' : k) :
  Region u v → Region u' v' → sol u v = sol u' v' → u = u' ∧ v = v' := by
    intro h₁ h₂ h₃; simp_all +decide [ sol ] ;
\end{lstlisting}

\end{document}
